\documentclass[exercice]{thedocument}%
\startdatakeys
\source{originale}
\cycle{4}
\competencesDuSocle{RA, CA}
\connaissancesRequises{A3a*}
\competencesTravaillees{A3b4,A3c*,A3d1}
\enddatakeys
\begin{document}%
Soient trois angles $\widehat{\text{A}}$, $\widehat{\text{B}}$ et $\widehat{\text{C}}$.
On sait que~:
\begin{itemize}
    \item $\widehat{\text{A}}=\widehat{\text{B}}$ ;
    \item $\widehat{\text{C}}=50^\circ$ ;
    \item Les trois angles forment ensemble un angle droit.
\end{itemize}
\begin{enumerate}
    \item Traduire le troisième point ci-dessus par une égalité mathématiques.
    \item On note $x$ la mesure des angles $\widehat{\text{A}}$ et  $\widehat{\text{B}}$.
    Déduire de la question précédente une équation d'inconnue $x$.
    \item En résolvant l'équation, en déduire la valeur des angles $\widehat{\text{A}}$ et $\widehat{\text{B}}$.
\end{enumerate}
\end{document}
